\section{Methods}\label{sec:methods}
\subsection{Evaluation price and pre-registration}\label{sec:methods-prereg}
% src: S-PAPER-2, S-PAPER-3
Each market has a scheduled close time $T$.
We set $T$ to the \texttt{close\_at} field.
Trades can continue after the close, often once the answer is known.
A price taken after the close can then contain the answer.
Then calibration looks better than it is.
In a one-month pilot of \PilotNMarkets{} markets, \PilotShareMarketsAfterClose{} had a trade after \texttt{close\_at}.
Those late trades are \PilotShareVolumeAfterClose{} of the pilot volume.
So we use the scheduled close and not the last trade.

\Cref{fig:timeline} shows where the evaluation price sits.
The evaluation price \pOne{} is a mean price.
It averages the trades at the last timestamp at or before $T-\DaysShort$ day.
We drop a market if this price is older than \MaxStaleHours{} hours.
We also drop a market with fewer than \MinTrades{} trades before that time.
A price from a nearly empty market is not a crowd view.
We report Q1 with and without this filter.

\begin{figure}[H]
\centering
\begin{tikzpicture}[font=\small]
\draw[thick,-{Stealth}] (0,0) -- (12.2,0);
\draw[thick] (0.2,-0.12) -- (0.2,0.12) node[above] {market opens};
\draw[thick,blue!60!black] (7.0,-0.2) -- (7.0,0.2) node[above] {$T-\DaysShort$ day};
\draw[thick] (9.2,-0.12) -- (9.2,0.12) node[above] {$T$ = \texttt{close\_at}};
\foreach \x in {10.0,10.5,11.0,11.5} \draw[red!70!black,thick] (\x,-0.15) -- (\x,0.15);
\node[below=6pt,blue!60!black,align=center] at (7.0,0) {evaluation price \pOne};
\node[below=6pt,red!70!black,align=center] at (10.75,0) {trades after the close};
\end{tikzpicture}
\caption{Where the evaluation price sits in a market's life (illustrative). Why: a price taken after the close can already contain the answer. Source: pre-registration, \texttt{analysis/prereg\_mvp1.md}.}
\label{fig:timeline}
\end{figure}

The criteria were fixed before any Q1 result.
\Cref{tab:prereg} lists them (\Cref{sec:app-prereg}).
The verdict applies to the full-data run with the same parameters.

% src: S-PAPER-3
One rule changed after we saw a full-data Q1 result.
The original rule was the last price at or before $T-\DaysShort$ day.
Many trades share one block timestamp, so this rule left their order undefined.
The new rule is the mean price of the trades at the last timestamp.
We made the change for determinism, not because of the result.
It came after we saw a full-data Q1 result, so we report it as an amendment.
The amendment in the pre-registration file records that the verdict did not change.
All numbers in this paper come from the runs with the amended rule.
This paper uses an effect size and an interval that differ from the pre-registration: per-bin gaps and the z interval.
It drops the Gini coefficient named in our proposal.
Both changes keep the analysis within the lecture contents.
The verdict rule and the verdict sentence are unchanged.
These methods are outside the lectures: the Lorenz curve, Wilson intervals, the cluster bootstrap chi-square, logistic calibration, and the Brier skill score.
All other methods come from the CS 6344 course lectures.

\subsection{Q1: calibration}\label{sec:methods-q1}
% src: S-PAPER-2, S-PAPER-3
We treat each outcome as a Bernoulli trial with success probability equal to the price.
Outcome 1 is the first-listed outcome of a market, labeled Yes in \OutcomeOneYesShare{} of kept markets.
So the outcome-1 count in a bin is approximately binomial.
We form \NBins{} equal-width price bins, and we use equal-frequency bins as a check.
In bin $b$, $n_b$ is the market count, $\bar p_b$ the mean price, and $O_b$ the outcome-1 count.
The null hypothesis is that the outcome-1 rate equals the mean price.
The effect size is the per-bin gap, the observed rate minus the mean price.
The expected outcome-1 count is $E_b=n_b\bar p_b$.
We test each bin with a binomial test.
We give a \CiLevel{} z interval for the observed rate $\hat r_b=O_b/n_b$, which is $\hat r_b\pm1.96\sqrt{\hat r_b(1-\hat r_b)/n_b}$.
We correct the \NBins{} tests with Bonferroni and with the Benjamini-Hochberg (BH) procedure \citep{benjamini1995controlling}.
The overall test is the goodness-of-fit form of the chi-square test: observed against expected counts, with the expected counts taken from the prices.
It runs on the bin by \{outcome 1, outcome 0\} table:
\begin{equation}
\chi^2=\sum_{b}\left[\frac{(O_b-E_b)^2}{E_b}+\frac{(O_b-E_b)^2}{n_b-E_b}\right].
\label{eq:chisq}
\end{equation}
We also report the mean squared residual of the price, $\frac1n\sum(p-y)^2$, where $n$ is the total market count.
The pre-registration calls it the Brier score \citep{brier1950verification}.
We add the log loss and the mean entropy $H(p)$ in bits.
We treat the market as a classifier that says outcome 1 wins when $p>\FavoriteLine$.
We compare the price distributions of markets where outcome 1 won and markets where outcome 0 won with a Kolmogorov-Smirnov (KS) test.
The KS statistic $D$ is the effect size.
For the longshot check we look at the two cheapest bins.

\subsection{Q2: wallet concentration}\label{sec:methods-q2}
% src: S-PAPER-2, S-PAPER-10
We work with wallet volume in USDC, by role.
Volume is the USDC sum over every trade in the kept markets.
It includes trades after $T-\DaysShort$ day.
We did not measure trade-count shares.
We report the share of role volume held by the top \TopOneLabel, \TopFiveLabel, \TopTenLabel, and \TopTwentyLabel{} of wallets.
We compare the top \TopTwentyLabel{} share with the 80/20 reference.
This is descriptive, with no pre-registered verdict line.
We draw the complementary CDF (CCDF) of wallet volume on log-log axes.
We sample its ranks from the top, so the largest wallets get points.
We also draw the wallet distribution, the volume distribution, and the Lorenz curve of each role.
Each is split at the top \TopTwentyLabel{} of wallets.
The Lorenz curve is outside the lectures.
We use it only to show how the top shares build up.
As a pipeline check, we compare our top \TopOneLabel{} maker share with the dataset paper's.
We test the log volume of makers against takers with a Welch t-test and a KS test.
The effect size is the difference of the mean log10 volume.
We report Spearman's $\rho$ between trade count and volume.

% src: S-PAPER-2, S-PAPER-6
As a side analysis, we run k-means on wallet features (\Cref{sec:app-kmeans}).

\subsection{Q3: distributions}\label{sec:methods-q3}
% src: S-PAPER-2, S-PAPER-5
We summarise price, trade size, and wallet volume with their mean and median.
We plot histograms of the price \pOne{} and of trade size.
We draw QQ plots of trade size and of wallet volume, before and after a log transform.
We fit a normal, an exponential, and a power law to trade size and wallet volume.
The power law is fitted on the top \TopTenLabel{} of values only, above a lower bound $x_{\min}$.
We measure each fit with the KS statistic $D$ against the data.
% src: S-PAPER-13
This follows the power-law method of \citet{clauset2009power}.
Unlike that method, we fix $x_{\min}$ at the top \TopTenLabel{} and do not choose it by the smallest $D$.
\todo{owner: why $x_{\min}$ is fixed at the top \TopTenLabel{} instead of chosen by the smallest $D$}
Trade sizes come from a deterministic sample of \NTradeSizeSample{} trades from the kept markets.

\subsection{Reproducibility}\label{sec:methods-repro}
% src: S-PAPER-5, S-PAPER-6
All data queries ran on the UTRGV Cradle cluster (SLURM, CPU only).
The first-stage run is \RunId{} (job \SlurmJobId), with code version \CodeVersion.
The software and code hashes are in \Cref{sec:app-repro}.
